\chapter{Professional Money Managers and their Influence}

These notes follow Robert J. Shiller's Lecture 20 on professional money managers and institutional investors, curated by LazyingArt LLC. The lecture begins with a question that sounds almost bureaucratic and turns out not to be bureaucratic at all: who actually controls wealth, influence, and decision-making in a modern economy? Shiller answers that question by giving us perspective before taxonomy. We first look at what households own, then at what they owe, then at what is missing even from net worth, and only after that do we turn to fiduciary duty, regulation, funds, trusts, pensions, endowments, family offices, and foundations.

\section{Institutional Investors in Perspective}

The opening move matters. ``Institutional'' sounds dull, but the lecturer insists that we are really talking about the people who control much of the wealth of the world. They are largely unseen; nobody makes movies about them. Yet they increasingly determine what gets financed, which risks are borne, and what arrangements ordinary households can rely on. In that sense, institutional investment is part of the governance of the world.

So we begin, as the lecture does, by looking not at a legal definition but at the balance sheet of U.S. households and nonprofits. Shiller takes his numbers from Table B-100 of the Federal Reserve's balance-sheet accounts and asks us first to look at the total. Measured household assets come to about \$70 trillion. That is not everything valuable in the nation, but it is the stock of assets the Federal Reserve can place on a balance sheet.

A few large categories already show the structure of the story:
\begin{align}
\text{Real estate} &\approx 18 \text{ trillion}, &
\text{Pension funds} &\approx 13 \text{ trillion}, \\
\text{Non-corporate business equity} &\approx 6 \text{ trillion}, &
\text{Deposits} &\approx 8 \text{ trillion}, \\
\text{Corporate equities} &\approx 8 \text{ trillion}, &
\text{Mutual funds} &\approx 5 \text{ trillion}.
\end{align}

The point of this list is not that every item is institutionally held. Real estate and non-corporate business are much more directly tied to households and families. But pension funds, bank deposits, mutual funds, life insurance reserves, and much of the bond market are already managed through professional institutions. By the lecturer's estimate, less than half but close to half of measured household assets are already in institutional hands. A century earlier, almost none of this would have been true. Society is becoming more professionalized, more institutionalized, and more dependent on specialized money managers.

Still, Shiller is careful not to erase the family. The family remains economically important. The lecture will keep that fact in view from beginning to end.

\section{Assets, Liabilities, and Household Net Worth}

At this point the lecture pauses and asks the obvious corrective question: what about liabilities? Gross assets are not net wealth. If households own assets and also owe money, we have to subtract.

\begin{figure}[htbp]
\centering
\includegraphics[width=0.84\textwidth]{lecture_20_figure_02.png}
\caption{Household liabilities and net worth for U.S. households and nonprofits, 2010-IV.}
\end{figure}

Let us write measured household assets as \(A_{\mathrm{hh}}\), measured household liabilities as \(L_{\mathrm{hh}}\), and net worth as \(NW_{\mathrm{hh}}\). The slide gives the core accounting identity directly:
\begin{align}
NW_{\mathrm{hh}}
  &= A_{\mathrm{hh}} - L_{\mathrm{hh}} \\
  &= 70{,}740 - 13{,}918 \\
  &= 56{,}822,
\end{align}
where the units are billions of dollars.

The liability structure on the slide is simple but revealing:
\begin{align}
\text{Home mortgages} &= 10{,}070, \\
\text{Consumer credit} &= 2{,}435, \\
\text{Loans \& other} &= 1{,}413, \\
\text{Total liabilities} &= 13{,}918.
\end{align}

The biggest household liability is mortgage debt. In the spoken lecture Shiller gives the more conversational version of the same point: households have about \$15 trillion of real estate and about \$10 trillion of mortgage debt, so the net position in housing is much smaller than the gross real-estate number suggests. The spoken figures are rounded and do not reconcile exactly line by line, but the economic lesson is plain: one must net the debt against the asset.

He then translates net worth into per-capita terms. The slide gives about \$184{,}000 per person. For a family of four,
\begin{equation}
4 \times 184{,}000 = 736{,}000,
\end{equation}
which the lecturer rounds upward conversationally to ``almost \$800{,}000.'' The rhetoric is deliberate: on average the country looks astonishingly wealthy. But the caveat is immediate. It is only an average. Wealth is very unequally distributed.

Shiller then makes one more adjustment, not as a formal national-accounting identity, but as a matter of perspective. Households will ultimately bear public debt through taxes, so he mentally charges federal debt and state-local debt back against household wealth. Using the slide values,
\begin{equation}
56.822 - 14.286 - 3.135 \approx 39.401 \text{ trillion dollars}.
\end{equation}
In lecture style, this becomes ``something like \$40 trillion or less.''

That adjustment is intentionally informal. The point is not to rewrite the Federal Reserve's sectoral accounts. The point is to discipline our intuition. Once private liabilities and public obligations are taken seriously, measured household wealth is still large, but it is smaller than the raw asset total suggests.

And now comes the next widening move. Even after we have subtracted liabilities, have we really counted the wealth of the nation?

\section{Human Capital and the Scale of National Wealth}

Shiller answers that question by undercutting his own topic. Institutional investors are important, but if we want complete perspective we must include something that does not appear on the household balance sheet at all. The transcript briefly says ``human rights'' at this point, but the lecture and the blackboard make clear that he means human capital.

Human capital is the value of people and what they can do and produce. If total national wealth is meant seriously, it must include the present value of future human income.

\begin{figure}[htbp]
\centering
\includegraphics[width=0.84\textwidth]{lecture_20_figure_03.png}
\caption{Human-capital wealth calculated by capitalizing national income.}
\end{figure}

The board organizes the idea in two steps: inputs on the left, valuation on the right. The lecture's input block is
\begin{align}
Y_0 &= 13 \text{ trillion dollars per year}, \\
g &= 0.03, \\
r &= 0.05,
\end{align}
where \(Y_0\) is current national income, \(g\) is the growth rate, and \(r\) is the discount rate.

A cautious standard reconstruction of the board's calculation is the Gordon-growth present-value formula
\begin{equation}
W = PV = \frac{Y_0}{r-g}.
\end{equation}

Substituting the lecture's parameters gives
\begin{equation}
W = \frac{13\ \text{trillion}}{0.05-0.03}.
\end{equation}

The visible board then records the lecture's own comparison number:
\begin{equation}
W = \frac{13\ \text{trillion}}{0.05-0.03} = 260\ \text{trillion}.
\end{equation}

\begin{remark}
The algebraic move on the board is clear, but the numerical arithmetic is rough if read literally. The safe reading is that the lecture is using a Gordon-style capitalization to establish scale, not to present a fully audited national-wealth estimate. What matters is the order-of-magnitude point: once human capital is counted, total national wealth is far larger than the measured household balance sheet.
\end{remark}

Shiller then repeats the move at the world level. Taking world GDP, about \$62 trillion, as a proxy for world income, we may write the same kind of expression,
\begin{equation}
W_{\mathrm{world}} \sim \frac{62\ \text{trillion}}{r-g}.
\end{equation}
The lecture gives the resulting world wealth as roughly \$1.2 quadrillion. Again, the exact arithmetic is less important than the perspective-setting purpose. The world economy is enormous, growing, and increasingly organized through professional financial institutions, but institutionally managed wealth is still only a fraction of all wealth if we count productive human capacity.

\subsection{Question \& Answer}

\paragraph{Question.} If institutional investors manage so much, why are they still only a fraction of total wealth?

\paragraph{Answer.} Because the balance sheet of tradable or measurable assets is not the whole economy. Human capital does not sit on the household asset list, but it is still wealth in the economically serious sense: it is the present value of future productive capacity. Once we capitalize national income, total wealth becomes much larger than the stock of assets directly managed by institutions. Shiller therefore deliberately ``diminishes'' his own lecture. Institutional investors are important, but they do not exhaust national wealth, and they do not erase the family as an economic unit.

\section{From Family Provision to Institutional Management}

This widening of perspective leads directly to the next question. If the family remains so important, why does institutional management keep growing?

The lecture's answer is social and practical rather than ideological. More and more, modern societies solve old family problems through formal institutions. Retirement is the clearest example. In an older family economy, aging parents move in with their children and are supported directly. In the institutional version, they save through pension funds, then choose among housing and care arrangements, including assisted living facilities. The lecturer's emphasis is on choice. Institutionalization is not simply a loss of warmth; it is also an enlargement of possible lives.

The same logic appears in healthcare. Families once tried to provide care directly, and often badly. Modern health plans and benefit structures are more institutional and, for all their defects, more competent. The same is true of risk management. A family, even a clever family, is too small a unit to diversify risk well. Professional managers can pool risks across many households and increasingly across countries.

That does not make the system infallible. Shiller points directly to the recent financial crisis and the failure to manage mortgage risk properly. Professionalization comes with growing pains. Risk managers are becoming more important, but sometimes they fail exactly where they are supposed to excel. That is why the lecture now turns to fiduciary duty and the law of prudence. Once we hand our savings and our risks to professionals, how are they supposed to behave?

\section{Fiduciary Duty, Prudence, and Post-Crisis Regulation}

Shiller's next sentence in effect is: now they have a fiduciary duty. If one is managing other people's money, the law does not allow the attitude, ``It is not my money; what do I care?''

\begin{definition}
A fiduciary is a person or institution charged with managing resources on behalf of others and therefore bound to act in the interest of the beneficiary rather than merely in personal interest.
\end{definition}

The lecture then introduces the prudent person rule. ERISA in 1974 gives the important U.S. formulation, and the core idea is that a fiduciary should manage with the care, skill, prudence, and diligence that a prudent person familiar with such matters would use in a like enterprise and with like aims.

That sounds reasonable at once. A pension manager should not speculate in wild schemes or racehorses. But the lecture immediately asks the destabilizing question: what is a prudent person? The difficulty is that prudence can become social imitation. A fiduciary may feel compelled not to do what seems truly intelligent, but to do what other respectable fiduciaries are expected to do. The rule can therefore suppress unconventional but defensible judgment.

The lecture gives a strong historical example. For much of the twentieth century, university endowments sat heavily in government bonds because that looked prudent. Then portfolios such as Yale's, under David Swensen and others, adopted a broader and more intellectually ambitious interpretation of prudence. That loosened the old convention. It also helped make institutional investors more aggressive, and Shiller's point is that this changing interpretation was one ingredient in the broader bubble that led to the crisis. The lecture is subtle here. It is not saying that all aggression is bad. It is saying that ``prudence'' is unstable when it is defined by convention alone.

\subsection{Question \& Answer}

\paragraph{Question.} What does it mean to invest prudently when convention itself may be wrong?

\paragraph{Answer.} It means that prudence cannot simply be whatever respectable people happen to be doing. A prudent-person rule begins as a restraint on recklessness, but it can become a demand for conformity. In one period that conformity forces institutions into overcautious portfolios; in another it loosens and admits excessive leverage or speculative enthusiasm. Shiller's point is that prudence is not self-interpreting. Law can state the duty, but it cannot fully remove the problem of judgment.

That unresolved problem is what drives the next move, from fiduciary language to post-crisis regulatory language.

\begin{figure}[htbp]
\centering
\includegraphics[width=0.84\textwidth]{lecture_20_figure_04.png}
\caption{Dodd-Frank (2010) introduced against earlier prudent-person and ERISA notes.}
\end{figure}

Dodd-Frank is presented as the most important U.S. financial legislation since the Great Depression. Shiller reports a revealing textual fact: the phrase ``prudent person'' appears nowhere in the Act, while ``prudential standards'' appears repeatedly. The emphasis shifts from a broad fiduciary norm toward explicit regulatory supervision of institutional risk-taking.

Leverage supplies the concrete example. In the lecture, leverage is borrowing to buy assets, and therefore a measure of the risk imposed on a portfolio by debt-financed exposure. Before the crisis the economy became increasingly leveraged. Dodd-Frank responds by creating the FSOC, the Financial Stability Oversight Council, to recommend leverage standards and broader prudential standards for financial firms and to place them under stronger regulatory authority.

The historical drift is therefore striking. Power moved from individual investors toward institutional investors, and after the crisis some of that power shifts again toward the state. The same pattern, Shiller suggests, is visible internationally and also in the declining trust placed in private rating agencies such as Moody's and Standard \& Poor's.

\section{Advisors, Funds, and Trust Structures}

At this point Shiller broadens the picture. Institutional finance is not only a matter of giant pools of capital. It also reaches households through the people and structures that stand between ordinary savers and complex markets.

He says explicitly that he wants to talk about financial advisors. Advisors do not necessarily manage the portfolio directly. Rather, they give advice to people who do not know how to manage it themselves. That alone makes them socially important. In the U.S. setting of the lecture, financial advisors are regulated, effectively through SEC registration and FINRA-administered approval and education. Shiller also mentions NAPFA, the National Association of Personal Financial Advisors, as a professional organization. In the lecture's concrete description, such advisors may charge roughly \$75 to \$300 per hour and operate under codes of conduct and loyalty to the client.

Financial planners occupy a murkier position. The transcript is noisy at this point, but the institutional point is clear enough: the boundary between advice, planning, and regulated responsibility is not perfectly clean. Dodd-Frank calls for further study of planners, and the mortgage-broker example shows why this matters. Before the crisis, mortgage brokers were much less tightly licensed than one might have expected, even though households increasingly depended on them to navigate borrowing decisions. The larger problem is that ordinary life is now more financial and less directly familial, while the advice markets that mediate this shift are still imperfect.

From there the lecture turns to some central institutional forms. The mutual fund is one of them. It is an investment company owned by its participants, investing typically in stocks and passing returns through to them. Shiller recalls the Massachusetts Investment Trust of the 1920s as an early model of an honest public fund: open portfolio reporting, no privileged senior class, and a transparent division of returns. After the abuses of the 1920s, the Investment Company Act of 1940 set the legal stage for the growth of mutual funds in the United States.

Europe develops an analogous vehicle in UCITS, the Undertakings for Collective Investment in Transferable Securities. The lecture's emphasis is not doctrinal detail but functional comparison. Europe standardized what had previously been fragmented national arrangements. One technical difference Shiller stresses is the treatment of capital gains: in the U.S. mutual-fund form, gains realized within the fund can be passed through to investors even if a given investor has not sold; under the UCITS structure, taxation is more closely tied to the investor's own sale. For that reason the lecture suggests that the European form has some practical appeal.

The trust is the next great institutional invention in the lecture. Here the tone changes. Trusts are not merely funds for the mass public; they are devices for making one person's intention survive time.

\begin{definition}
A trust is a legal arrangement in which assets are held and managed on behalf of another person or for a specified purpose. A spendthrift trust is a trust that pays income to the beneficiary while restricting access to the principal.
\end{definition}

The lecture's classic example is a child who cannot manage his or her own affairs and may outlive the parents. A trust allows the parent to create an income stream that continues after the parent's death, whether the trustee is a trust company or a person who can appoint a successor. The legal separation matters: if the trustee later goes bankrupt, the trust assets are not simply available to the trustee's creditors.

The spendthrift trust sharpens the point. Parents who fear that a child may dissipate an inheritance can restrict the child to income while protecting principal. The lecture gives a further motive that is characteristic of Shiller's institutional imagination: protection in divorce. The point is not moralism. It is that trust law creates durable structures for handling incapacity, improvidence, family conflict, and long horizons.

\section{Pensions, Endowments, and the Problem of Delegated Judgment}

The lecture then turns to pensions, which are perhaps the clearest instance of finance solving a basic human problem. People get old. They cannot keep working forever. Historically their fate depended heavily on the fortunes and duties of children, employers, or local charity. The pension plan is one of the institutional inventions by which modern societies changed that condition.

Shiller gives the history in episodes. The first U.S. pension plan, he says, was American Express in 1875. This was not yet the credit-card firm of modern memory but a delivery company. Its promise was already recognizably ``defined benefit'': after sufficient service, older disabled employees would receive half the average pay of their last ten years for life. In compact notation we may write the model promise as
\begin{equation}
B = 0.5 \times \big(\text{average of the last 10 years' pay}\big).
\end{equation}

That formula captures the central managerial problem. If the benefit is fixed in advance, then someone must decide today how much money to set aside and how to invest it so that the future liability can actually be met. Pension management therefore becomes an asset-liability problem, not merely a return-seeking problem.

After American Express came large industrial and union plans, but the early history is full of failure. The crash after 1929 wiped out many promises. Union pension funds often performed disastrously. The world, Shiller says, was amateurish in this domain.

The General Motors pension plan of 1950 marks a decisive improvement because it was fully funded. The company did not merely promise future income; it set aside and invested assets now. The lecture's analogy is elegant: this is like a trust for workers. If the company itself later dies, the workers' pension assets do not have to die with it.

Studebaker, by contrast, becomes the warning case. Its pension plan was inadequately funded, and bankruptcy exposed the weakness of labor-management bargaining as a guarantee of long-run retirement security. That episode leads directly to ERISA, the Employee Retirement Income Security Act of 1974, and to the PBGC, the Pension Benefit Guaranty Corporation, which insures pension funding and forces plans into greater scrutiny.

The lecture then makes the next historical turn. ERISA was written for a world dominated by defined benefit plans. But beginning in the 1980s firms increasingly shifted toward defined contribution plans. The reason is managerial difficulty. If the employer promises the retirement benefit itself, the manager must hit a moving target: future wages, future returns, future longevity, future required funding. If instead the employer promises only the contribution, then the contractual burden shifts.

\subsection{Question \& Answer}

\paragraph{Question.} Why did firms move from defined benefit promises to defined contribution plans?

\paragraph{Answer.} Because a defined benefit plan fixes the retirement liability in advance, while a defined contribution plan fixes only the current payment into the worker's account. In the first case the manager must solve a hard target-hitting problem; in the second case the final retirement income is whatever the accumulated portfolio can support. The shift therefore eases the burden on the firm and the manager, but it also moves more responsibility onto the household.

Schematically, if we let \(C\) denote the employer's current contribution, then the defined contribution structure says, in effect: the promise is \(C\), not \(B\). The worker later retires on the realized value of the portfolio rather than on a fixed pension formula.

Shiller immediately adds that this new system has its own growing pains. In defined contribution plans workers often choose whether to participate, how much to contribute, and how to allocate assets across equity, bonds, real estate, or other funds. That gives freedom, but it also produces nonparticipation, overconcentration in risky assets, and what he calls something like amateur investing. Recent reforms therefore make it easier for firms to provide advice, automatically enroll employees, and place silent participants into default ``prudent'' allocations.

From pensions the lecture moves naturally to endowments. Endowments are again institutional portfolios serving a purpose beyond themselves. And here the history is equally instructive. Yale, in 1825, placed its entire endowment in the Eagle Bank of New Haven and lost it all. Boston University later concentrated a large endowment position in a single genetic-engineering company and suffered heavily. The University of Bridgeport exhausted its endowment and had to alter its institutional path drastically. These are not stray anecdotes. They are pedagogical illustrations of how amateur or concentrated stewardship can destroy a mission.

Against that background, David Swensen appears not merely as a successful investor but as a manager of delegated judgment. His task is to preserve the university's ability to pursue its educational purpose. That is why the lecture treats endowment management as morally and institutionally serious. The payoff is concrete: students whose tuition and living support are financed by successful endowment strategy are living on the institutional consequences of prudent, but not merely conventional, financial management.

\section{Family Offices, Family Foundations, and Philanthropic Capital}

The lecture closes by returning to the family. This is not an afterthought. It is the resolution of a tension that has been present since the opening. Institutionalization has advanced enormously, but the family has not disappeared. It survives and acts through institutional shells.

Shiller says this explicitly. Marx predicted the end of the family as an economic unit. It did not happen. The family is too deeply rooted in social life. So he comes back to family offices and family foundations.

A family office is a private institutional arrangement for managing the wealth of a single wealthy family. The lecture gives the scale intuition in the simplest possible arithmetic:
\begin{equation}
0.05 \times 100\ \text{million dollars} = 5\ \text{million dollars per year}.
\end{equation}
If a family has \$100 million in assets and earns about five percent, it has enough annual income to justify dedicated staff. That is why a family office may employ two, three, or five people full-time to manage the portfolio and plan related structures such as trusts for children.

A family foundation is different. It is not chiefly for the family's private portfolio but for a charitable purpose carried by the family name. The lecture suggests that many wealthy families have both: a family office for wealth management and a family foundation for philanthropy.

\begin{figure}[htbp]
\centering
\includegraphics[width=0.84\textwidth]{lecture_20_figure_05.png}
\caption{Family offices, family foundations, and the board's 36{,}000 count.}
\end{figure}

The board records the striking number \(36{,}000\), and the lecture identifies this as the number of family foundations in the United States as of 2006. The chalkboard itself shows only the count; the time and place qualification come from the spoken lecture. Shiller's point is not merely statistical. It is that increasing affluence and changing values are pushing more families to organize charitable action in durable institutional form.

The lecture also democratizes the idea. A typical family foundation need not be huge. It may begin with one or two million dollars, endowed to some named cause, perhaps neighborhood improvement or another public mission. The tax treatment matters as well: charitable deductions encourage the irrevocable transfer of funds into the foundation, allowing wealth to accumulate there while the donor thinks seriously about long-run purpose rather than responding impulsively to every solicitation.

This is also one more way in which institutions solve family problems rather than replace the family. A foundation can outlive its founder, give children an ongoing shared role, and organize a family's identity around a cause rather than around consumption alone.

That prepares the last moral example. Drawing on Robert Frank's \emph{Richistan}, Shiller gives the two-sided case of Paul Allen. On one side there is the 400-foot yacht, the spectacle of private excess. On the other side there is the Paul G. Allen Family Foundation and the fact of more than a billion dollars already donated to charity. The point is not to canonize or condemn. It is to make us reserve judgment and notice that great wealth often arrives institutionally twice: once as display and once as philanthropy.

The lecture ends there for a reason. Institutional investing is not only about power, leverage, and regulation. It is also about the forms through which private wealth can be made to serve durable public purposes.

\section{Summary}

The lecture unfolds as a sequence of enlargements and corrections. We begin with the claim that institutional investors are part of the governance of the modern world. We then inventory U.S. household assets and see that professional institutions already manage a striking share of measured wealth. We subtract liabilities and discover that gross assets are not net worth. We then widen the frame again and find that even net worth is not total wealth unless human capital is included.

Only after that perspective-setting work does Shiller turn to social function, fiduciary duty, prudence, Dodd-Frank, advisors, funds, trusts, pensions, endowments, and family structures. The through-line is consistent throughout. Modern finance is not merely an industry category. It is a set of institutional inventions for handling retirement, care, risk sharing, intergenerational support, and philanthropy. These inventions fail, they are reformed, and they become more professional. But the family does not disappear inside this process. It remains, increasingly, as an actor that works through offices, trusts, foundations, and other institutional forms.